\chapter{بهبودهای کلاسیک و انتخاب مرکز}

% \section{الگوریتم هارتیگان--وونگ}

% روش هارتیگان--وونگ \cite{hartigan1979} برخلاف لوید فقط به دنبال پایداری
% نسبت به انتساب هم‌زمان نزدیک‌ترین مرکز نیست؛ پس از مقداردهی اولیه، انتقال
% تک‌نقطه‌ای میان خوشه‌ها را مستقیماً از نظر کاهش مجموع مربعات درون‌خوشه‌ای
% می‌سنجد. برای نقطه‌ی $x$ در خوشه‌ی مبدأ $L_1$ با اندازه‌ی $n_1$، هزینه‌ی
% حذف آن متناسب با
% \begin{equation}
% \frac{n_1}{n_1-1}\lVert x-\mu_{L_1}\rVert^2
% \end{equation}
% است و افزودن آن به خوشه‌ی مقصد $L$ با اندازه‌ی $n_L$ هزینه‌ی
% \begin{equation}
% \frac{n_L}{n_L+1}\lVert x-\mu_L\rVert^2
% \end{equation}
% دارد. انتقال زمانی انجام می‌شود که کاهش خالص مثبت باشد. نقطه‌ای که تنها
% عضو خوشه‌ی خود است منتقل نمی‌شود تا خوشه‌ی تهی ایجاد نشود.

% پیاده‌سازی مرجع دو مرحله را متناوب می‌کند:
% \begin{enumerate}
% \item در \lr{OPTRA}، مقصدهای ممکن برای یک نقطه بررسی و انتقال بهبوددهنده
% انجام می‌شود. مفهوم «مجموعه‌ی زنده» از بررسی دوباره‌ی خوشه‌هایی که از آخرین
% تغییرشان ثابت مانده‌اند جلوگیری می‌کند.
% \item در \lr{QTRAN}، فقط بهترین مقصد جایگزینِ ذخیره‌شده برای هر نقطه سریع
% بررسی می‌شود. اگر انتقال انجام شود، مراکز و اندازه‌ها بلافاصله اصلاح می‌شوند.
% \end{enumerate}
% فرایند با نبود انتقال بهبوددهنده خاتمه می‌یابد. در نتیجه، جواب نسبت به
% انتقال یک نقطه یک بهینه‌ی محلی قوی‌تر است، اما همچنان تضمین بهینه‌ی سراسری
% ندارد. فایل \lr{interactive/hartigan-wong-vs-lloyd.html} تفاوت انتقال‌های
% تک‌نقطه‌ای با دورهای دسته‌ای لوید را به‌صورت مفهومی و گام‌به‌گام نمایش می‌دهد.

\section{مقداردهی اولیه‌ی \lr{k-means++}}

آرتور و واسیلیویتسکی نشان دادند که انتخاب دقیق‌تر مراکز اولیه می‌تواند
کیفیت مورد انتظار \lr{k-means} را کنترل کند \cite{arthur2007}. الگوریتم
\lr{k-means++} چنین است:
\begin{enumerate}
    \item مرکز نخست به‌صورت یکنواخت از نقاط داده انتخاب می‌شود.
    \item اگر $D(x)$ فاصله‌ی نقطه‌ی $x$ تا نزدیک‌ترین مرکز انتخاب‌شده باشد،
          مرکز بعدی با احتمال
          \begin{equation}
              P(x)=
              \frac{D(x)^2}{\sum_{x'\in\mathcal{X}}D(x')^2}
          \end{equation}
          انتخاب می‌شود.
    \item گام دوم تا انتخاب $k$ مرکز تکرار و سپس الگوریتم معمول لوید اجرا می‌شود.
\end{enumerate}
این نمونه‌گیری «وزن‌دهی $D^2$» نام دارد. قضیه‌ی اصلی مقاله بیان می‌کند
\begin{equation}
    \mathbb{E}[\Phi]\leq 8(\ln k+2)\Phi_{\mathrm{OPT}}.
    \label{eq:kpp-bound}
\end{equation}
بنابراین روش در بدترین حالت
$O(\log k)$-رقابتی است و برای نمونه‌های خوب می‌تواند کران ثابت
داشته باشد. روش پیشینِ استروفسکی و همکاران، کران ثابت را تحت شرط
\begin{equation}
    \frac{\Phi_{\mathrm{OPT},k}}
    {\Phi_{\mathrm{OPT},k-1}}\leq\varepsilon^2
\end{equation}
به دست می‌داد؛ یعنی جواب $k$-خوشه‌ای باید به‌طور محسوسی بهتر از جواب
$(k-1)$-خوشه‌ای باشد.

\begin{lemma}[هویت مرکز ثقل]
    برای مجموعه‌ی متناهی $S$، مرکز ثقل $c(S)$ و هر نقطه‌ی $z$ داریم
    \begin{equation}
        \sum_{x\in S}\lVert x-z\rVert^2
        =
        \sum_{x\in S}\lVert x-c(S)\rVert^2
        +|S|\lVert c(S)-z\rVert^2.
        \label{eq:centroid-identity}
    \end{equation}
\end{lemma}

تفسیر رابطه‌ی \ref{eq:centroid-identity} این است که هزینه‌ی انتخاب هر
نماینده‌ی دلخواه برابر است با هزینه‌ی (فاصله‌ی) نقاط نسبت به مرکز ثقل واقعی،
به‌علاوه‌ی جریمه‌ی دور بودن آن نماینده از مرکز ثقل که در تعداد نقاط ضرب
می‌شود. 
\section{سایر مسیرهای مطالعه‌شده}

\subsection{روش \lr{global k-means}}

روش \lr{global k-means} \cite{likas2003} مراکز را یکی‌یکی اضافه می‌کند. برای
افزایش تعداد خوشه‌ها از $k-1$ به $k$، هر نقطه‌ی داده یک‌بار به‌عنوان مرکز
جدید در کنار جواب قبلی آزمایش و لوید اجرا می‌شود؛ بهترین جواب نگه داشته
می‌شود. این جست‌وجوی قطعی حساسیت به مقداردهی اولیه را کاهش می‌دهد، اما
اجرای مکرر آن پرهزینه است. نسخه‌ی سریع مقاله با یک کران کاهش خطا، نامزدها
را محدود می‌کند.

\subsection{\lr{X-means}}

\lr{X-means} \cite{pelleg2000} از یک کران پایین برای $k$ آغاز می‌کند و
خوشه‌ها را به‌صورت محلی برای شکستن به دو بخش نامزد می‌کند. این فرایند تا کران بالای $k$ ادامه
می‌یابد. مزیت اصلی آن جست‌وجوی هم‌زمان ساختار و تعداد خوشه‌هاست، ولی نتیجه
به مدل آماری و بازه‌ی مجاز $k$ وابسته است.
